\documentclass[11pt]{amsart}
\usepackage{amsmath,amssymb,amsthm}
\usepackage[margin=1.05in]{geometry}

\numberwithin{equation}{section}
\newtheorem{theorem}{Theorem}[section]
\newtheorem{proposition}[theorem]{Proposition}
\newtheorem{lemma}[theorem]{Lemma}
\newtheorem{corollary}[theorem]{Corollary}
\theoremstyle{definition}
\newtheorem{definition}[theorem]{Definition}
\theoremstyle{remark}
\newtheorem{remark}[theorem]{Remark}
\newtheorem{example}[theorem]{Example}

\newcommand{\Z}{\mathbb{Z}}
\newcommand{\Q}{\mathbb{Q}}
\newcommand{\R}{\mathbb{R}}
\newcommand{\N}{\mathbb{N}}
\newcommand{\C}{\mathbb{C}}
\newcommand{\T}{\mathbb{T}}
\newcommand{\cA}{\mathcal{A}}
\newcommand{\cG}{\mathcal{G}}
\newcommand{\cK}{\mathcal{K}}
\newcommand{\cO}{\mathcal{O}}
\newcommand{\hO}{\widehat{\cO}}
\newcommand{\cT}{\mathcal{T}}
\newcommand{\Vect}{\mathrm{Vec}}
\newcommand{\pp}{\mathfrak{p}}
\newcommand{\qq}{\mathfrak{q}}
\newcommand{\aaa}{\mathfrak{a}}
\newcommand{\bb}{\mathfrak{b}}
\newcommand{\Ns}{\mathrm{N}}
\newcommand{\Prob}{\operatorname{Prob}}
\newcommand{\Idl}{\mathbb{I}}
\newcommand{\Cl}{\operatorname{Cl}}
\newcommand{\Iso}{\operatorname{Iso}}
\newcommand{\Lat}{\operatorname{Lat}}
\newcommand{\coker}{\operatorname{coker}}
\newcommand{\Ind}{\operatorname{Ind}}
\newcommand{\cR}{\mathcal{R}}
\newcommand{\Ent}{\mathrm{S}}
\newcommand{\Msym}{M_{\mathrm{sym}}}
\newcommand{\phisym}{\varphi_{\mathrm{sym}}}
\newcommand{\spn}{\operatorname{span}}
\newcommand{\id}{\mathrm{id}}

\PassOptionsToPackage{hyphens}{url}
\usepackage[hidelinks,bookmarks=false]{hyperref}

\title[Ideals, boundary, and index]{Localized Bost--Connes systems II:\\ ideal structure, boundary quotient, and index}
\author{Ruqing Chen}
\address{GUT Geoservice Inc., Montreal, Canada}
\email{ruqing@hotmail.com}
\thanks{DOI: \href{https://doi.org/10.5281/zenodo.22177726}{10.5281/zenodo.22177726}.
Sources, code and data:\newline\url{https://github.com/Ruqing1963/localized-bost-connes}
(this paper: \texttt{papers/II-structure/}).}
\date{}

\begin{document}

\begin{abstract}
We study the ideal structure, the boundary quotient, the $K$-theory and the index theory
of the localized Bost--Connes system $\cA_{K,S}=C(Y_S)\rtimes J_S$, and show that all of it
is governed by one structural feature: since $J_S$ is free abelian on $S$, each prime
carries exactly \emph{one} isometry and the range projections
$\mathbf 1_{\aaa Y_S}$ are \emph{nested} rather than orthogonal.

From this we obtain: the groupoid is topologically principal, its isotropy on the finite
closed orbit $\partial Y_S\cong\Cl^+_S$ is $I_S^0=\ker(I_S\to\Cl^+_S)$, and off
$\partial Y_S$ it is trivial whenever $K$ has finitely many totally positive units; the ideal
lattice is then the lattice of open invariant subsets of $\widetilde Y_S\setminus\partial Y_S$
glued to the lattice of closed subsets of the torus $\widehat{I_S^0}$, the latter
contributing ideals that neither the diagonal nor the gauge action sees. The tight quotient
makes the isometries unitary rather than Cuntz, giving
$\partial\cA_{K,S}\cong M_{h^+_S}(C^*(I_S^0))$, a type $\mathrm I$ algebra with a faithful
trace whose only isomorphism invariant is the narrow class number $h^+_S$, recoverable as
the divisibility of $[1]$ in $K_0=\Lambda^{\mathrm{ev}}(I_S^0)$. For a single prime,
$K_1(\cA_{K,S})=0$ and $K_0(\cA_{K,S})\cong C(G_S,\Z)/\Z\oplus\Z$ is free; for several primes
$K_*$ is the abutment of a collapsing Koszul spectral sequence which we do not evaluate.

The same feature governs the index theory. At the symmetric $\mathrm{KMS}_\beta$ state, every
prime inclusion $M_{S_1}\subset M_{S_2}$ has infinite index for \emph{every} conditional
expectation, for the elementary reason that the valuation projections $\mathbf 1_{\{v_\pp=k\}}$
lie in the relative commutant and are moved onto one another by partial isometries commuting
with $M_{S_1}$; there is no index formula in $\Ns\pp$ and $[\Cl^+_{S_2}:\Cl^+_{S_1}]$, and no
Galois correspondence with intermediate prime sets, there being a continuum of intermediate
algebras where there are two intermediate prime sets. The finite index instead comes from
symmetry breaking: $M_\varepsilon\cong\Msym^{\widehat{\Xi_\beta}}\subset\Msym$ is the trivial
inclusion of a factor tensored with $\C\subset\ell^\infty(\widehat{\Xi_\beta})$, of index
$|\Xi_\beta|$, its intermediate algebras are indexed by the partitions of
$\widehat{\Xi_\beta}$, and $\log[\Msym:M_\varepsilon]=|A_\beta|\log2$ is also the relative
entropy of an extremal phase against the symmetric one.

Finally we run the control experiment. Restoring the additive part gives
$\cA^{\mathrm{aff}}_{K,S}=C(\hO_S)\rtimes(\cO_K\rtimes J_S)$, whose additive cosets provide
\emph{finite covers} and hence genuine Cuntz relations; it is a simple purely infinite
Kirchberg algebra. Nested versus partitioning range projections is therefore
exactly the dividing line. Even there, however, $[1]$ carries almost nothing: the Cuntz
relation forces its order to divide $\gcd_{\pp\in S}(\Ns\pp-1)$, which is $1$ as soon as
$S$ contains a prime of norm $2$.
\end{abstract}

\maketitle

\tableofcontents

\section{The single-isometry phenomenon}

We keep the notation of \cite{Main}: $K$ is a number field, $S$ a set of finite primes,
$J_S=\bigoplus_{\pp\in S}\N$ the free abelian monoid of integral ideals supported on $S$
with enveloping group $I_S$, $Y_S=\hO_S\times_{U_S}G_S$, and
\[
\cA_{K,S}=C(Y_S)\rtimes J_S=\overline{\spn}\bigl\{\mu_\aaa f\mu_\bb^*\bigr\},
\qquad \mu_\aaa^*\mu_\aaa=1,\quad \mu_\aaa\mu_\aaa^*=\mathbf 1_{\aaa Y_S}.
\]
\emph{The series.} The foundational paper \cite{Main} classifies the $\mathrm{KMS}$ states
of $\cA_{K,S}$ and introduces the obstruction group $\Xi_\beta$. Five papers build on it,
numbered I--V:
\[
\begin{array}{ll}
\text{I} & \text{factor type and the spectrum of transitions}\\
\text{II} & \text{ideal structure, boundary quotient, and index\quad(this paper)}\\
\text{III} & \text{invariants, entropy, and the measure class}\\
\text{IV} & \text{noncommutative metric geometry}\\
\text{V} & \text{the free variant and the free critical exponent}
\end{array}
\]
Paper I has appeared; Papers III--V are in preparation. \emph{Nothing in this paper depends
on them}: every statement below is proved here or cited from \cite{Main}, and the forward
references are signposts only. We refer to them by numeral, without citation, for that
reason.

Recall the exact sequence \cite[(2.1)]{Main}
\begin{equation}\label{eq:exact}
1\to U_S/\overline{\cO^\times_{K,+}}\to G_S\to\Cl^+_S\to1,\qquad h^+_S:=|\Cl^+_S| .
\end{equation}

\begin{lemma}[Single isometry, nested projections]\label{lem:single}
$J_S$ is free abelian on $S$. Hence:
\begin{enumerate}
\item for each $\pp\in S$ there is exactly one generating isometry $\mu_\pp$;
\item the range projections are \emph{nested}:
$\mathbf 1_{\aaa Y_S}\mathbf 1_{\bb Y_S}=\mathbf 1_{\mathrm{lcm}(\aaa,\bb)Y_S}$;
\item there is no finite family of isometries in $\cA_{K,S}$ with mutually orthogonal ranges
summing to $1$; indeed
$Y_S\setminus\bigcup_{\pp\in S}\pp Y_S=\{[x,g]:x\in\hO_S^\times\}$ is nonempty and clopen.
\end{enumerate}
\end{lemma}

\begin{proof}
Freeness is unique factorization of ideals. (2) is immediate from
$\aaa Y_S\cap\bb Y_S=\mathrm{lcm}(\aaa,\bb)Y_S$. For (3), two of the range projections
$\mathbf 1_{\aaa Y_S}$ are never orthogonal by (2); and if $x$ is a unit at every $\pp$ then
$[x,g]\notin\pp Y_S$ for any $\pp$, which gives the last assertion. That no family of
isometries whatever has orthogonal ranges summing to $1$ --- $\cA_{K,S}$ is not properly
infinite --- follows from Theorem~\ref{thm:boundary} below: $\cA_{K,S}$ maps unitally onto
an algebra with a faithful tracial state $\tau$, and $n$ isometries with orthogonal ranges
summing to $1$ would give $n=\sum\tau(v_i^*v_i)=\sum\tau(v_iv_i^*)=1$.
\end{proof}

Lemma~\ref{lem:single} is the source of every structural statement below. We flag its use
by referring to it explicitly. The contrast is with the $ax+b$ algebras of Cuntz--Li and
Cuntz--Deninger--Laca, where each ideal $\aaa$ carries $\Ns\aaa$ isometries indexed by the
additive cosets $a+\aaa$, whose ranges \emph{partition}; \S\ref{sec:affine} restores that
ingredient and confirms that it is the cause.

\section{Isotropy, topological principality, and the ideal lattice}

Let $\cG_S=I_S\ltimes\widetilde Y_S$ be the dilated groupoid of \cite[Prop.~2.8]{Main},
with points $[x,g]$ subject to $(wx,g)\sim(x,r(w)g)$ for $w\in U_S$.

\begin{lemma}[Isotropy formula]\label{lem:iso}
For $y=[x,g]$ put $T(y)=\{\pp\in S:x_\pp=0\}$. Then
\begin{equation}\label{eq:iso}
\Iso_{\cG_S}(y)=\bigl\{\aaa\in I_{T(y)}:\ \sigma_\aaa\in r(U_{T(y)})\bigr\}.
\end{equation}
\end{lemma}

\begin{proof}
$\aaa\cdot[x,g]=[x,g]$ means $\aaa x=wx$ and $g=r(w)^{-1}\sigma_\aaa g$ for some $w\in U_S$.
Comparing $\pp$-components: if $x_\pp\ne0$ then $\pp^{a_\pp}x_\pp=w_\pp x_\pp$ forces
$a_\pp=0$ and $w_\pp=1$; if $x_\pp=0$ both sides vanish. So $\aaa\in I_T$, $w\in U_T$, and
the second equation is $\sigma_\aaa=r(w)$.
\end{proof}

\begin{proposition}[Isotropy off the boundary]\label{prop:nojump}
Write $\Idl_T=\{u\in\Idl_S:\ u_{\pp'}=1\ \text{for all }\pp'\in S\setminus T\}$. Then for
$y$ with $T(y)=T$,
\begin{equation}\label{eq:isodiv}
\Iso_{\cG_S}(y)\ =\ \mathrm{div}\Bigl(\overline{K^\times_{S,+}}\cap\Idl_T\Bigr).
\end{equation}
If $\cO^\times_{K,+}$ is finite --- in particular if $K=\Q$ or $K$ is imaginary quadratic
--- then for every proper $T\subsetneq S$ the right-hand side is trivial, so
$\Iso_{\cG_S}(y)=\{1\}$: the action of $I_S$ on $\widetilde Y_S\setminus\partial Y_S$ is free.

At the bottom stratum the isotropy is nontrivial:
\[
\partial Y_S:=\bigcap_{\aaa\in J_S}\aaa Y_S=\bigl\{[0,g]\bigr\}\cong G_S/r(U_S)\cong\Cl^+_S,
\]
a set of $h^+_S$ points, on which $I_S$ acts transitively through the class map with
isotropy $I_S^0:=\ker(I_S\to\Cl^+_S)$, of index $h^+_S$.
\end{proposition}

\begin{proof}
By Lemma~\ref{lem:iso}, $\aaa\in I_T$ lies in the isotropy iff $\sigma_\aaa\in r(U_T)$, i.e.\
iff some idele $s$ with $\mathrm{div}(s)=\aaa$ satisfies $s\in U_T\cdot\overline{K^\times_{S,+}}$.
Writing $s=wu$ with $w\in U_T$ and $u\in\overline{K^\times_{S,+}}$, the element $u=w^{-1}s$
is trivial at every $\pp'\notin T$ and has $\mathrm{div}(u)=\aaa$, which is \eqref{eq:isodiv};
the converse is immediate.

Suppose now $\cO^\times_{K,+}$ is finite. Then $K^\times_{S,+}$ is discrete in $\Idl_S$: an
element of $K^\times_{S,+}$ lying in the neighbourhood $\prod_{\pp\in F}(1+\pp^k\cO_\pp)\times\prod_{\pp\notin F}\cO_\pp^\times$
of $1$ has trivial divisor, hence is a totally positive unit, and is $1$ for $k$ large. So
$\overline{K^\times_{S,+}}=K^\times_{S,+}$, and $u=\alpha\in K^\times_{S,+}$ is trivial at some
$\pp'\in S\setminus T$. Under the diagonal embedding the $\pp'$-component of $\alpha$ is
$\alpha$ itself, and $K\to K_{\pp'}$ is injective, so $\alpha=1$ and $\aaa=(1)$.

For the bottom stratum, $\bigcap_\aaa\aaa Y_S$ forces $x_\pp=0$ for every $\pp$, so $x=0$;
the fibre is $G_S/r(U_S)\cong\Cl^+_S$ by \eqref{eq:exact}, and \eqref{eq:iso} with $T=S$
becomes $\{\aaa:[\aaa]=1\}=I_S^0$.
\end{proof}

\begin{remark}[The general case, and why it does not matter here]\label{rem:leopoldt}
When $\cO^\times_{K,+}$ is infinite the right-hand side of \eqref{eq:isodiv} can in
principle be nontrivial: one needs $\alpha^{-1}$ to lie in the closure of
$\cO^\times_{K,+}$ inside $\prod_{\pp'\in S\setminus T}\cO^\times_{\pp'}$, a condition of
Leopoldt type. We do not decide it. Nothing below depends on it:
Theorem~\ref{thm:principal} needs only that the strata are nowhere dense, which holds
whatever their isotropy, since for infinite $S$ every basic clopen set contains points with
$x_\pp\ne0$ at any prescribed $\pp$.
\end{remark}

\begin{remark}\label{rem:whynojump}
The natural expectation is that killing the coordinates in $T$ frees the corresponding
translations, giving isotropy $I_T$. It frees them in the $\hO_S$-direction but not in the
$G_S$-direction: multiplication by $\aaa$ translates $g$ by $\sigma_\aaa$, whose components
at the primes \emph{outside} $T$ are nontrivial. Over $\Q$ with $T=\{2\}$ and $\aaa=(2)$,
the components of $\sigma_\aaa$ at $3,5,7$ are $2,2,2$, none equal to $1$.
\end{remark}

\begin{theorem}[Topological principality, and the two halves of the ideal lattice]\label{thm:principal}
Let $\partial I:=C^*(\cG_S|_{\widetilde Y_S\setminus\partial Y_S})\cap\cA_{K,S}$ be the ideal of
$\cA_{K,S}$ corresponding to the open invariant set $\widetilde Y_S\setminus\partial Y_S$, so
that $\cA_{K,S}/\partial I=\partial\cA_{K,S}$ is the boundary quotient of
Theorem~\ref{thm:boundary}.
\begin{enumerate}
\item $\cG_S$ is topologically principal, for every $K$ and $S$; consequently $C(Y_S)$ is a
Cartan subalgebra of $\cA_{K,S}$ \cite{Renault08}.
\item The ideals of $\cA_{K,S}$ containing $\partial I$ correspond bijectively to the closed
subsets of the torus $\widehat{I_S^0}$, via
$\partial\cA_{K,S}\cong M_{h^+_S}(C(\widehat{I_S^0}))$. These ideals are not detected by the
diagonal $C(Y_S)$: two of them have the same intersection $C_0(Y_S\setminus\partial Y_S)$
with it. They are not all gauge-invariant either: the gauge action is the translation flow
of $\widehat{I_S^0}$ by the characters $\aaa\mapsto\Ns\aaa^{it}$, which is nontrivial, so
a closed subset that is not a union of flow orbits --- a point, for instance --- gives an
ideal that the gauge action moves.
\item If the action of $I_S$ on $\widetilde Y_S\setminus\partial Y_S$ is free --- by
Proposition~\ref{prop:nojump}, whenever $\cO^\times_{K,+}$ is finite --- then the ideals of
$\partial I$ correspond bijectively and order-isomorphically to the open $I_S$-invariant
subsets of $\widetilde Y_S\setminus\partial Y_S$, and the diagonal detects them.
\item Consequently, in that case,
$\operatorname{Prim}(\cA_{K,S})=(\widetilde Y_S\setminus\partial Y_S)/I_S\ \sqcup\ \widehat{I_S^0}$
as a set, the first piece open and the second closed.
\end{enumerate}
\end{theorem}

\begin{proof}
(1) By \eqref{eq:iso} every unit with nontrivial isotropy lies in some stratum
$\{x_\pp=0\ \forall\pp\in T\}$ with $T\ne\emptyset$, hence in some $\{x_\pp=0\}$, a closed set
with empty interior since $\cO_\pp$ has no isolated points. A countable union of such sets
is meagre, so the units with trivial isotropy are dense. (This does not use
Proposition~\ref{prop:nojump}.) The restriction $\cG_S|_{Y_S}$ is then topologically
principal as well, and $\cA_{K,S}=C^*(\cG_S|_{Y_S})$ with diagonal $C(Y_S)$, which is Cartan
by Renault's theorem \cite{Renault08}.

(2) Ideals containing $\partial I$ are ideals of the quotient, which is $M_h(C(X))$ with
$X=\widehat{I_S^0}$ by Theorem~\ref{thm:boundary}; its ideals are the $M_h(C_0(X\setminus Z))$,
$Z\subseteq X$ closed. The image of $C(Y_S)$ in $M_h(C(X))$ is the algebra
$C(\partial Y_S)=\C^h$ of diagonal matrices with constant entries, which meets
$M_h(C_0(X\setminus Z))$ trivially for every nonempty $Z$; so every ideal $J\supseteq\partial I$
with $J\ne\cA_{K,S}$ has $J\cap C(Y_S)=\partial I\cap C(Y_S)=C_0(Y_S\setminus\partial Y_S)$.
The gauge action $\sigma_t(\mu_\aaa)=\Ns\aaa^{it}\mu_\aaa$ descends to the unitaries
$\mu_\aaa$, $\aaa\in I_S^0$, of $C^*(I_S^0)=C(X)$, where it is translation by the character
$\chi_t:\aaa\mapsto\Ns\aaa^{it}$; $\chi_t\ne1$ for $t\ne0$ since $\log\Ns\aaa\ne0$ for
$\aaa\ne(1)$, so the flow moves every point of $X$.

(3) $\cG_S|_{\widetilde Y_S\setminus\partial Y_S}$ is second countable, \'etale, amenable and
principal, so by Renault \cite{Renault} its $C^*$-algebra has ideals in bijection with the
open invariant subsets, via $V\mapsto C^*(\cG_S|_V)$, and its diagonal is a Cartan
subalgebra detecting them; $\partial I$ is the full corner $\mathbf 1_{Y_S}(\cdot)\mathbf 1_{Y_S}$,
which preserves the lattice. (4) is (2) and (3) together with the exactness of
$0\to C^*(\cG_S|_U)\to C^*(\cG_S)\to C^*(\cG_S|_{U^c})\to0$ for amenable \'etale groupoids.
\end{proof}

\begin{remark}\label{rem:lattice}
Part (2) shows that ``the ideal lattice is the lattice of open invariant sets'' is false
for $\cA_{K,S}$: the boundary orbit has isotropy $I_S^0$ of infinite rank, and its
$C^*$-algebra $C(\widehat{I_S^0})$ contributes a torus of primitive ideals invisible to the
diagonal. Topological principality of $\cG_S$ gives the correspondence only on the
essentially principal part $\widetilde Y_S\setminus\partial Y_S$. When $\cO^\times_{K,+}$ is
infinite, Remark~\ref{rem:leopoldt} leaves open whether some proper stratum has nontrivial
isotropy; if it does, that stratum contributes further ideals of the same kind.
\end{remark}

\begin{proposition}[The lattice is not the power set of $S$]\label{prop:lattice}
Let $K=\Q$, $S=\{\ell\}$, so $Y_S=\Z_\ell$ and $\cA=C(\Z_\ell)\rtimes\N$. Then $\cA$ is a
full corner in $C_0(\Q_\ell)\rtimes\Z$ and there is an exact sequence
\begin{equation}\label{eq:oneprime}
0\longrightarrow C(\Z_\ell^\times)\otimes\cK\longrightarrow C_0(\Q_\ell)\rtimes\Z\longrightarrow C(\T)\longrightarrow0 .
\end{equation}
The ideal lattice of $\cA$ contains an order-isomorphic copy of the lattice of open subsets
of the Cantor set $\Z_\ell^\times$, hence is uncountable.
\end{proposition}

\begin{proof}
By Laca's dilation theorem $\cA$ is a full corner in $C_0(\Q_\ell)\rtimes\Z$, the dilation
being the limit of $\Z_\ell\xleftarrow{\ell}\Z_\ell\xleftarrow{\ell}\cdots$. On
$\Q_\ell\setminus\{0\}$ the action of $\Z$ by multiplication by $\ell$ is free and proper
with quotient $\Z_\ell^\times$, giving the ideal in \eqref{eq:oneprime}; the fixed point
$\{0\}$ gives the quotient $C^*(\Z)=C(\T)$. Ideals of $C(\Z_\ell^\times)\otimes\cK$
correspond to open subsets of $\Z_\ell^\times$ and pass to $\cA$ by the corner
correspondence.
\end{proof}

\begin{remark}
The description ``hereditary subsets of $\mathcal P(S)$'' would be correct for a graph or
Cuntz--Krieger algebra. Here the stratification by vanishing coordinates is present, but
each stratum retains a Cantor set of unit orbits, each closed set of orbits contributes an
ideal, and the quotient $C(\T)$ in \eqref{eq:oneprime} is the $h^+_S=1$, $I_S^0=\Z$ case of
Theorem~\ref{thm:principal}(2).
\end{remark}

\section{The boundary quotient}

\begin{theorem}[The tight quotient is abelian-by-finite]\label{thm:boundary}
Let $\partial\cA_{K,S}$ be the quotient by the ideal corresponding to
$\widetilde Y_S\setminus\partial Y_S$, equivalently the tight quotient in the sense of Exel
\cite{Exel}, equivalently the universal quotient in which every $\mu_\pp$ becomes unitary.
Then
\[
\partial\cA_{K,S}\ \cong\ C(\Cl^+_S)\rtimes I_S\ \cong\ M_{h^+_S}\bigl(C^*(I_S^0)\bigr)
\ \cong\ M_{h^+_S}\bigl(C(\widehat{I_S^0})\bigr).
\]
It is unital, separable, nuclear, satisfies the UCT, is of type $\mathrm I$ and carries a
faithful tracial state. It is not simple when $I_S^0\ne0$, and is never purely infinite; in
particular it is not a Kirchberg algebra.
\end{theorem}

\begin{proof}
Restricting the groupoid to the closed invariant set $\partial Y_S$ gives
$C^*(I_S\ltimes\partial Y_S)=C(\Cl^+_S)\rtimes I_S$ by Proposition~\ref{prop:nojump}. In the
quotient $\mathbf 1_{\pp Y_S}$ restricts to $1$ on $\partial Y_S\subseteq\pp Y_S$, so
$\mu_\pp$ is unitary; conversely imposing unitarity kills exactly the functions vanishing on
$\partial Y_S$. For a discrete group $\Gamma$ with a subgroup $\Gamma_0$ of finite index $h$,
Green's imprimitivity theorem \cite{Green} gives $C(\Gamma/\Gamma_0)\rtimes\Gamma\cong M_h(C^*(\Gamma_0))$,
and $I_S^0$ is free abelian so $C^*(I_S^0)$ is commutative. The remaining properties are
those of $M_h(C(X))$; a unital algebra with a tracial state is never purely infinite.
\end{proof}

\begin{remark}[Why no Cuntz relation is available]\label{rem:nocuntz}
By Lemma~\ref{lem:single}(2) the projections $\mathbf 1_{\aaa Y_S}$ are nested. Imposing
$\sum_\aaa\mu_\aaa\mu_\aaa^*=1$ over any finite family would force all but one summand to
vanish. The tight relation therefore degenerates to $\mathbf 1_{\pp Y_S}=1$, which is
unitarity, and this is what Theorem~\ref{thm:boundary} implements. The Cuntz alternative
requires a partition; see \S\ref{sec:affine}.
\end{remark}

\begin{theorem}[The class of the unit records $h^+_S$]\label{thm:Kboundary}
$K_0(\partial\cA_{K,S})\cong\Lambda^{\mathrm{ev}}(I_S^0)$ and
$K_1(\partial\cA_{K,S})\cong\Lambda^{\mathrm{odd}}(I_S^0)$, and under this identification
\begin{equation}\label{eq:unit}
\bigl[1_{\partial\cA_{K,S}}\bigr]=h^+_S\cdot e_0 ,
\end{equation}
$e_0$ the generator of $\Lambda^0=\Z$. Hence $h^+_S$ is the largest integer dividing $[1]$,
and is an isomorphism invariant of the pair $(\partial\cA_{K,S},[1])$.
\end{theorem}

\begin{proof}
$K$-theory is Morita invariant and $C^*(\Gamma)\cong C(\widehat\Gamma)$ for $\Gamma$ free
abelian has $K_0=\Lambda^{\mathrm{ev}}(\Gamma)$, $K_1=\Lambda^{\mathrm{odd}}(\Gamma)$ with
$[1]$ the generator of $\Lambda^0$. The isomorphism $K_0(M_h(B))\cong K_0(B)$ sends
$[1_{M_h(B)}]$ to $h[1_B]$, since $1_{M_h(B)}$ is a sum of $h$ orthogonal projections each
equivalent to $1_B$; and $\Lambda^{\mathrm{ev}}$ is free abelian, so $h\,e_0$ is divisible by
exactly $h$.
\end{proof}

\begin{theorem}[Rigidity of the boundary]\label{thm:rigid}
For infinite $S_1\subseteq P_{K_1}$ and $S_2\subseteq P_{K_2}$,
\[
\partial\cA_{K_1,S_1}\cong\partial\cA_{K_2,S_2}\iff h^+_{S_1}=h^+_{S_2} .
\]
The boundary determines the order of the subgroup of the narrow class group generated by
$S$, and nothing else: not $K$, not $S$, not $\Cl^+_S$ as a group.
\end{theorem}

\begin{proof}
($\Leftarrow$) For infinite $S$, $I_S$ is free abelian of countably infinite rank and
$I_S^0$ has finite index, hence the same rank; so the two algebras are
$M_h(C(X))$ with the same $h$ and homeomorphic $X$.
($\Rightarrow$) Every irreducible representation of $M_h(C(X))$ has dimension exactly $h$,
and isomorphisms preserve irreducible representations and their dimensions; alternatively
apply Theorem~\ref{thm:Kboundary}. For the last statement, $C(\Cl^+_S)\rtimes I_S$ depends
on $\Cl^+_S$ only through the index of $I_S^0$.
\end{proof}

\section{$K$-theory of $\cA_{K,S}$}

The Stacey crossed product imposes the covariance relation $\mu_\pp\mu_\pp^*=\alpha_\pp(1)$,
which is Katsura covariance; hence for a single prime $\cA_{K,S}$ is the Cuntz--Pimsner
algebra $\cO_X$ of the associated correspondence, \emph{not} the Toeplitz algebra $\cT_X$.
Pimsner's $KK$-equivalence $\iota:A\to\cT_X$ therefore does not apply, and one must use his
exact sequence.

\begin{lemma}\label{lem:cantorspace}
$Y_S$ is compact, metrizable and totally disconnected without isolated points, hence a
Cantor set. Consequently $K^0(Y_S)=C(Y_S,\Z)$ is free abelian of countably infinite rank and
$K^1(Y_S)=0$.
\end{lemma}

\begin{proof}
$\hO_S$ is profinite, $G_S$ is a quotient of the totally disconnected $\Idl_S$ by a closed
subgroup hence profinite, and $Y_S$ is their quotient by the compact group $U_S$. A point
$[x,g]$ is not isolated: every neighbourhood contains $[x',g]$ with $x'\ne x$, since no
$\cO_\pp$ has isolated points.
\end{proof}

\begin{proposition}[One prime]\label{prop:K}
Let $S=\{\pp\}$ and let $\alpha_*$ be the endomorphism $\mathbf 1_U\mapsto\mathbf 1_{\pp U}$
of $K^0(Y_S)=C(Y_S,\Z)$. Then
\begin{equation}\label{eq:K}
K_1(\cA_{K,S})=\ker(1-\alpha_*)=0,
\qquad
K_0(\cA_{K,S})=\coker(1-\alpha_*)\ \cong\ C(G_S,\Z)\big/\Z\cdot1\ \oplus\ \Z ,
\end{equation}
a free abelian group. In particular $[1]=[\mathbf 1_{\pp Y_S}]$ in $K_0$, so
$\iota_*:K^0(Y_S)\to K_0(\cA_{K,S})$ is \emph{not} injective.
\end{proposition}

\begin{proof}
$\mu_\pp$ is an isometry with $\mu_\pp\mu_\pp^*=\mathbf 1_{\pp Y_S}$, so
$1\sim\mathbf 1_{\pp Y_S}$ and $[1]=[\mathbf 1_{\pp Y_S}]$. The six-term sequence for the
Cuntz--Pimsner algebra of a correspondence \cite[\S4]{Pimsner} (see \cite{Katsura} for the
covariance condition) with $K_1(C(Y_S))=0$ collapses to
$K_1(\cA)=\ker(1-\alpha_*)$, $K_0(\cA)=\coker(1-\alpha_*)$ on $C(Y_S,\Z)$.

If $f=\alpha_*f=(f\circ\pp^{-1})\mathbf 1_{\pp Y_S}$ then $f$ is supported in $\pp^nY_S$
for every $n$, i.e.\ in $\partial Y_S=\{[0,g]\}$, a closed set with empty interior; a
locally constant function supported there vanishes. So $K_1=0$.

For $K_0$ we use Laca's dilation \cite{Laca}: $\cA$ is a full corner in
$C_0(\widetilde Y_S)\rtimes\Z$, and by Pimsner--Voiculescu \cite{PV}
$K_0(C_0(\widetilde Y_S)\rtimes\Z)=H_0(\Z;C_c(\widetilde Y_S,\Z))$, the coinvariants,
$K_1=H_1$, the invariants. Write $\widetilde Y_S=\widetilde Y_\emptyset\sqcup\partial Y_S$
with $\widetilde Y_\emptyset=\{[x,g]:x\ne0\}$ open. The map
$[\pi^nu,g]\mapsto(n,r(u)g)$ identifies $\widetilde Y_\emptyset$ with $\Z\times G_S$, on which
$\pp$ acts by $(n,h)\mapsto(n+1,\sigma_\pp h)$, freely; hence
$H_1(\Z;C_c(\widetilde Y_\emptyset,\Z))=0$ and $H_0=C(G_S,\Z)$, the class of $f$ being
$\sum_n(\sigma_\pp^{-n})_*f(n,\cdot)$. On $\partial Y_S\cong\Cl^+_S$ the generator acts
through the class map, so $H_1=H_0=\Z$ by Shapiro's lemma. The long exact sequence of
$0\to C_c(\widetilde Y_\emptyset,\Z)\to C_c(\widetilde Y_S,\Z)\to C(\partial Y_S,\Z)\to0$
reads
\[
0\to H_1(\widetilde Y_S)\to\Z\xrightarrow{\ \delta\ }C(G_S,\Z)\to H_0(\widetilde Y_S)\to\Z\to0,
\]
and $\delta(1)$ is the class of $(1-\alpha_*)\mathbf 1_{Y_S}=\mathbf 1_{Y_S\setminus\pp Y_S}
=\mathbf 1_{\{0\}\times G_S}$, i.e.\ the constant function $1$; so $\delta$ is injective,
$H_1(\widetilde Y_S)=0$ again, and $H_0(\widetilde Y_S)$ is an extension of $\Z$ by
$C(G_S,\Z)/\Z\cdot1$, which splits. Both summands are free, $1$ being part of a basis of
$C(G_S,\Z)$.
\end{proof}

For $S$ finite with $|S|\ge2$ the primes interact: by \cite[Lem.~2.10]{Main} the generator
$\pp$ moves the $G_S$-coordinate by $\sigma_\pp$, whose components at the other primes are
nontrivial, so the action of $\Z^S$ on $\widetilde Y_S$ is not a product of actions on
the factors $K_\pp$. What survives is the following.

\begin{proposition}[Finite $S$]\label{prop:Kfinite}
For $S$ finite, $K_*(\cA_{K,S})=K_*(C_0(\widetilde Y_S)\rtimes\Z^S)$ is computed by the
spectral sequence $E^2_{p,q}=H_p(\Z^S;K_q(C_0(\widetilde Y_S)))\Rightarrow K_{p+q}$ of
Kasparov \cite[\S6]{Kasparov}. Since $K_1(C_0(\widetilde Y_S))=0$, only the row $q=0$ is
nonzero, the sequence collapses, and $K_0(\cA_{K,S})$ (resp.\ $K_1$) carries a finite
filtration whose subquotients are the even (resp.\ odd) Koszul homology groups
$H_p(\Z^S;C_c(\widetilde Y_S,\Z))$ of the commuting automorphisms $\alpha_{\pp*}$, $\pp\in S$.
The bottom piece $H_0(\Z^S;C_c(\widetilde Y_S,\Z))$, the coinvariants, is the image of
$K^0(Y_S)\to K_0(\cA_{K,S})$. For general $S$, $K_*(\cA_{K,S})=\varinjlim K_*(\cA_{K,F})$
over finite $F\subseteq S$, for any compatible system of embeddings as in
\cite[Prop.~2.7]{Main} (see Remark~\ref{rem:embed} below).
\end{proposition}

\begin{proof}
Laca's dilation, Kasparov's spectral sequence for actions of $\Z^n$, and the continuity of
$K$-theory under direct limits.
\end{proof}

We have not computed the Koszul homology for $|S|\ge2$, and we make no claim about
$K_*(\cA_{K,S})$ beyond Propositions \ref{prop:K} and \ref{prop:Kfinite}; the top group
$H_{|S|}$ is the group of $\Z^S$-invariant compactly supported functions, which vanishes by
the argument of Proposition~\ref{prop:K}, but the intermediate groups receive contributions
from the strata $\{x_\pp=0,\ \pp\in T\}$, on which $\Z^T$ acts on $G_S$ by translations
with nontrivial invariants.

\begin{example}[$K=\Q$, $S=\{\ell\}$]\label{ex:oneprime}
Here $G_S=\Z_\ell^\times$ and Proposition~\ref{prop:K} gives $K_1(\cA)=0$ and
$K_0(\cA)\cong C(\Z_\ell^\times,\Z)/\Z\oplus\Z\cong\Z^{(\infty)}$, in agreement with
\eqref{eq:oneprime}. A direct computation of the Smith normal form of $1-\alpha_*$ on
level-$n$ functions, for $\ell=2,3,5$ and $n\le3$, gives kernel $0$ and free cokernel of
rank $\ell^{n+1}-\ell^n$, with no torsion.
\end{example}

\begin{remark}[$K$-theory carries no arithmetic: the mechanism]\label{rem:noarith}
For a single prime, $K_0(\cA_{K,S})$ is free abelian and depends on $(K,\pp)$ only through
the profinite group $G_S$; for $K=\Q$ it is $\Z^{(\infty)}$ for every $\ell$. The reason
is structural rather than accidental: by Propositions \ref{prop:K} and \ref{prop:Kfinite}
the image of $K^0(Y_S)$ in $K_0(\cA_{K,S})$ is the group of coinvariants
$C(Y_S,\Z)/\sum_\pp(1-\alpha_{\pp*})C(Y_S,\Z)$, and a functional on $C(Y_S,\Z)$ descends
to it exactly when it is $\alpha_*$-invariant, whereas the arithmetically meaningful
functionals --- the $\mathrm{KMS}_\beta$ measures --- satisfy
$\omega\circ\alpha_{\pp*}=\Ns\pp^{-\beta}\omega$ with $\Ns\pp^{-\beta}\ne1$. The passage to
$K_0$ annihilates precisely the data of interest. This is taken up in Paper III of the
series.
\end{remark}

\section{Index theory}

\begin{remark}[Nested prime sets]\label{rem:embed}
For $S_1\subset S_2$ the embedding $\cA_{K,S_1}\hookrightarrow\cA_{K,S_2}$, $f\mapsto f\circ\mathrm{pr}$,
$\mu_\aaa\mapsto\mu_\aaa$, of \cite[Prop.~2.7]{Main} requires an $I_{S_1}$-equivariant
continuous surjection $\mathrm{pr}:Y_{S_2}\to Y_{S_1}$. By \eqref{eq:exact}, $Y_S$ is a
disjoint union over $\Cl^+_S$ of copies of $\hO_S/\overline{\cO^\times_{K,+}}$, and
$\Cl^+_{S_1}$ is a \emph{subgroup} of $\Cl^+_{S_2}$, not a quotient. The coordinate
projection $\hO_{S_2}\to\hO_{S_1}$ therefore descends canonically exactly when
$\Cl^+_{S_2}=\Cl^+_{S_1}$, i.e.\ when the primes of $S_2\setminus S_1$ have classes in
$\Cl^+_{S_1}$; this holds for $K=\Q$, for every $S_1$ generating the narrow class group,
and for all $S_1\subseteq S_2$ inside a fixed $S$ once $S_1$ is large enough. Otherwise
$\mathrm{pr}$ exists but involves the choice of a $\Cl^+_{S_1}$-equivariant retraction
$\Cl^+_{S_2}\to\Cl^+_{S_1}$, and the words ``$G_S=\varprojlim G_{S_k}$'' in
\cite[Prop.~2.7]{Main} must be read with this qualification (the uses made of that
proposition in \cite{Main} are over $\Q$). \emph{Throughout this section we assume
$\Cl^+_{S_2}=\Cl^+_{S_1}$}, so that $\mathrm{pr}$ is canonical.
\end{remark}

Let $S_1\subset S_2$, let $\varphi_2$ be a $\mathrm{KMS}_\beta$ state of $\cA_{K,S_2}$ with
$\zeta_{K,S_2}(\beta)=\infty$, and put $M_{S_i}=\pi_{\varphi_2}(\cA_{K,S_i})''$, where
$\cA_{K,S_1}\subseteq\cA_{K,S_2}$ via Remark~\ref{rem:embed}.

\begin{proposition}[The conditional expectation exists]\label{prop:E}
$\sigma^{\varphi_2}_t=\sigma_{-\beta t}$ leaves $\cA_{K,S_1}$ globally invariant, so by
Takesaki's theorem \cite{Takesaki} there is a unique faithful normal conditional expectation
$E:M_{S_2}\to M_{S_1}$ with $\varphi_2\circ E=\varphi_2$. If $\varphi_2$ is extremal, its
restriction to $\cA_{K,S_1}$ need not be, so $M_{S_1}$ need not be a factor.
\end{proposition}

\begin{proof}
$\sigma_t$ fixes $C(Y_{S_2})\supseteq C(Y_{S_1})$ pointwise and scales each $\mu_\aaa$, so
preserves $\cA_{K,S_1}$; the GNS vector of a KMS state is separating, so $\varphi_2$ is
faithful on $M_{S_2}$ with modular group $\sigma_{-\beta t}$. Extremality is not preserved by
restriction because the simplices are indexed by $\widehat{\Xi_\beta(S_2)}$ and
$\widehat{\Xi_\beta(S_1)}$ respectively.
\end{proof}

We write $\Ind(F)\in[1,\infty]$ for the Pimsner--Popa index of a faithful normal conditional
expectation $F:M_{S_2}\to M_{S_1}$, the inverse of the best constant $\lambda$ with
$F(x)\ge\lambda x$ for all $x\ge0$ \cite{PP}, and $[M_{S_2}:M_{S_1}]_0=\inf_F\Ind(F)$ for
the Kosaki--Longo index \cite{Kosaki,Longo}.

\begin{theorem}[Prime inclusions have infinite index]\label{thm:infinite}
Let $\pp\in S_2\setminus S_1$ and $e_k=\mathbf 1_{\{v_\pp=k\}}=\mathbf 1_{\pp^kY_{S_2}}-\mathbf 1_{\pp^{k+1}Y_{S_2}}$.
\begin{enumerate}
\item For every $\mathrm{KMS}_\beta$ state $\varphi_2$, the $e_k$ are pairwise orthogonal
projections in $M_{S_1}'\cap M_{S_2}$ with $\sum_ke_k=1$ and
$\varphi_2(e_k)=(1-\Ns\pp^{-\beta})\Ns\pp^{-k\beta}>0$; in particular $M_{S_1}'\cap M_{S_2}$
contains a copy of $\ell^\infty(\N)$.
\item Let $\varphi_2=\phisym$ be the symmetric $\mathrm{KMS}_\beta$ state, so that in the
coordinates of \cite[Lem.~2.10]{Main} its measure is $\nu=\pi_\beta\otimes\mathrm{Haar}$.
Then $\Ind(F)=\infty$ for every faithful normal conditional expectation
$F:M_{S_2}\to M_{S_1}$, hence $[M_{S_2}:M_{S_1}]_0=\infty$, already when a single prime is
added.
\end{enumerate}
\end{theorem}

\begin{proof}
(1) For $\qq\in S_1$, multiplication by $\qq$ does not change $v_\pp$, so
$\qq\{v_\pp\ge k\}=\{v_\pp\ge k\}\cap\qq Y_{S_2}$ and
\[
\mu_\qq\mathbf 1_{\{v_\pp\ge k\}}=\alpha_\qq\bigl(\mathbf 1_{\{v_\pp\ge k\}}\bigr)\mu_\qq
=\mathbf 1_{\{v_\pp\ge k\}}\mathbf 1_{\qq Y_{S_2}}\mu_\qq=\mathbf 1_{\{v_\pp\ge k\}}\mu_\qq ,
\]
using $\mathbf 1_{\qq Y}\mu_\qq=\mu_\qq\mu_\qq^*\mu_\qq=\mu_\qq$; taking adjoints gives
commutation with $\mu_\qq^*$, and commutation with $C(Y_{S_1})$ is automatic. The
$\mathbf 1_{\{v_\pp\ge k\}}$ therefore commute with $\cA_{K,S_1}$, hence with $M_{S_1}$, and so
do their differences $e_k$. The values $\varphi_2(e_k)$ come from
$\varphi_2(\mathbf 1_{\pp^kY})=\Ns\pp^{-k\beta}$ \cite[Prop.~2.8]{Main}, and
$\sum_ke_k=\mathbf 1_{\{v_\pp<\infty\}}=1$ in $M_{S_2}$ because $\varphi_2(v_\pp\ge k)\to0$.

(2) Let $\varphi_2=\phisym$. The GNS space is $H=L^2(\cR,\tilde\nu)$, $\cR$ the orbit
relation of $I_{S_2}$ on $(Y_{S_2},\nu)$ and $\tilde\nu$ its lift, $M_{S_2}=W^*(\cR,\nu)$,
and $e_k$ acts by multiplication by $\mathbf 1_{\{v_\pp=k\}}$ in the first variable
\cite[Prop.~6.1]{Main}. Fix $k,j\ge0$ and let $T:\{v_\pp=k\}\to\{v_\pp=j\}$,
$T(v,h)=(v+(j-k)e_\pp,h)$, in the coordinates $(v,h)\in V\times G_{S_2}$: $T$ moves the
$\pp$-valuation and nothing else. It is a homeomorphism between clopen sets with
$\nu\circ T=\Ns\pp^{-(j-k)\beta}\nu$ on its domain, since $\nu$ is a product measure with
geometric $\pp$-marginal and Haar $G_{S_2}$-marginal; and it commutes with the maps
$(v,h)\mapsto(v+e_\qq,\sigma_\qq h)$, $\qq\in S_1$, generating the action of $I_{S_1}$. Hence
$T$ maps $\cR$-classes to $\cR$-classes, and
\[
(W\xi)(y,y')=\Ns\pp^{(j-k)\beta/2}\,\xi(T^{-1}y,T^{-1}y'),\qquad y\in\{v_\pp=j\},
\]
is a partial isometry on $H$ with $W^*W=e_k$ and $WW^*=e_j$. Because $T$ commutes with the
$I_{S_1}$-action, does not move the coordinates seen by $C(Y_{S_1})\circ\mathrm{pr}$, and
has constant Radon--Nikodym derivative, $W$ commutes with $\pi(\cA_{K,S_1})$, hence
$W\in M_{S_1}'$ (the commutant in $B(H)$; $W$ need not lie in $M_{S_2}$).

Now let $F$ be a faithful normal conditional expectation with $\Ind(F)=\lambda^{-1}<\infty$,
so $F(e_k)\ge\lambda e_k$. Since $e_k$ commutes with $M_{S_1}$, $F(e_k)$ lies in
$Z(M_{S_1})$ and commutes with $e_k$; on the range of $e_k$ it is $\ge\lambda$, so
$e_k\le q_k:=\mathbf 1_{[\lambda,1]}\bigl(F(e_k)\bigr)\in M_{S_1}$, and
$\lambda q_k\le F(e_k)$ gives $\varphi_2(q_k)\le\lambda^{-1}\varphi_2(F(e_k))$. For $q\in M_{S_1}$
with $qe_k=e_k$ we get, using $W\in M_{S_1}'$,
\[
qe_j=qWW^*=WqW^*=W(qe_k)W^*=We_kW^*=e_j\qquad\text{for every }j,
\]
so $q=\sum_je_jq=1$. Applied to $q=q_k$: $1=\varphi_2(q_k)\le\lambda^{-1}\varphi_2(F(e_k))$,
i.e.\ $\varphi_2(F(e_k))\ge\lambda$ for every $k$. But $F(e_k)\in Z(M_{S_1})$ and
$\sum_kF(e_k)=1$, so $\sum_k\varphi_2(F(e_k))=1$, and $\varphi_2(F(e_k))\ge\lambda>0$ for
infinitely many $k$ is impossible.
\end{proof}

\begin{remark}[Scope]\label{rem:scope}
Part (1) holds for every KMS state; part (2) uses the product form of $\nu$ only to build
$W$. For a non-symmetric extremal state the fibre measure $\rho_v$ of \cite[Lem.~2.10]{Main}
is Haar measure on a coset that moves by $\sigma_\pp$ when $v_\pp$ increases, so $T$ no
longer preserves the measure class, and we have not decided whether every expectation has
infinite index in that case; the symmetric state is extremal, and $M_{S_2}$ a factor,
precisely when $\Xi_\beta(S_2)=\{1\}$, e.g.\ for $S_2=P_K$ and $\beta\le1$.
\end{remark}

\begin{remark}[Why no index formula can exist]\label{rem:noformula}
An index formula in $\Ns\pp$, a local unit index and $[\Cl^+_{S_2}:\Cl^+_{S_1}]$ would
combine finite quantities. It cannot exist because enlarging $S$ enlarges the \emph{base
space}: $Y_{S_2}$ carries a new $\pp$-adic coordinate, and the valuation $v_\pp$ alone,
before any unit or class-group data enters, already contributes an infinite-dimensional
commuting algebra. Those finite invariants do appear in this theory --- in
Theorem~\ref{thm:Kboundary} and in \eqref{eq:exact} --- but not in an index.
\end{remark}

\begin{proposition}[No Galois correspondence with prime sets]\label{prop:nogalois}
In the situation of Theorem~\ref{thm:infinite}(2) let $S_2=S_1\cup\{\pp\}$, so there are
exactly two intermediate prime sets. Then $\Lat(M_{S_1}\subset M_{S_2})$ has cardinality
at least $2^{\aleph_0}$.
\end{proposition}

\begin{proof}
For a partition $\mathcal P$ of $\N$ let $A_{\mathcal P}\subseteq W^*(e_k:k\ge0)\cong\ell^\infty(\N)$
be the subalgebra of functions constant on the blocks; it lies in the relative commutant, so
$N_{\mathcal P}=M_{S_1}\vee A_{\mathcal P}$ is intermediate. We recover $\mathcal P$ from
$N_{\mathcal P}$: an element $x=\sum_kc_ke_k$ of $W^*(e_k)$ lies in $N_{\mathcal P}$ iff for
every block $B$ the element $\sum_{k\in B}c_ke_k$ lies in $M_{S_1}\mathbf 1_B$, i.e.\ equals
$m\mathbf 1_B$ for some $m\in M_{S_1}$; conjugating by the partial isometries $W$ of the proof
of Theorem~\ref{thm:infinite} gives $c_je_j=me_j=W(me_k)W^*=c_ke_j$ for $j,k\in B$, so
$W^*(e_k)\cap N_{\mathcal P}=A_{\mathcal P}$. There are $2^{\aleph_0}$ partitions.
\end{proof}

We now locate the finite index. Fix $S$ and $\beta$ with $\Xi_\beta$ finite, and write
$G:=\widehat{\Xi_\beta}$, $\phisym$ for the $G_S$-invariant $\mathrm{KMS}_\beta$ state, and
$\varphi_\varepsilon$, $\varepsilon\in G$, for the extreme points.

\begin{lemma}[Central decomposition]\label{lem:central}
\begin{equation}\label{eq:central}
\Msym:=\pi_{\phisym}\bigl(\cA_{K,S}\bigr)''\ \cong\ \bigoplus_{\varepsilon\in G}M_\varepsilon,
\qquad M_\varepsilon:=\pi_{\varphi_\varepsilon}\bigl(\cA_{K,S}\bigr)'' ,
\end{equation}
with centre $Z(\Msym)\cong\C^{G}$ spanned by projections $z_\varepsilon$, each
$M_\varepsilon$ a factor, and $\omega_{\phisym}(z_\varepsilon)=|G|^{-1}$. The action
$\theta$ of $G_S$ extends to $\Msym$, factors through
$G_S/\Xi_\beta^\perp\cong G$, and permutes the $z_\varepsilon$ simply transitively.
\end{lemma}

\begin{proof}
By \cite[Thm.~3.6]{Main} the $\mathrm{KMS}_\beta$ simplex is $\Prob(G_S/\Xi_\beta^\perp)$,
so its extreme points are indexed by $G=\widehat{\Xi_\beta}$ and correspond to mutually
singular measures $\nu_\varepsilon$; hence the GNS representations are pairwise disjoint and
each $M_\varepsilon$ is a factor. The symmetric state is the uniform mixture
$\phisym=|G|^{-1}\sum_\varepsilon\varphi_\varepsilon$, so
$\pi_{\phisym}=\bigoplus_\varepsilon\pi_{\varphi_\varepsilon}$, giving \eqref{eq:central},
the centre, and $\omega_{\phisym}(z_\varepsilon)=|G|^{-1}$. By \cite[Thm.~3.6(3)]{Main} the
action of $G_S$ on the simplex factors through $G_S/\Xi_\beta^\perp$ and permutes the
extreme points simply transitively; it therefore permutes the $z_\varepsilon$ the same way.
\end{proof}

\begin{theorem}[The finite-index inclusion]\label{thm:finite}
Put $D:=M_{\varepsilon_0}$ for any $\varepsilon_0\in G$. The isomorphisms
$\theta_g:M_{\varepsilon_0}\to M_{g\varepsilon_0}$ identify the inclusion
$\Msym^{G}\subseteq\Msym$ with
\begin{equation}\label{eq:trivialincl}
D\otimes1\ \subseteq\ D\otimes\ell^\infty(G)=\bigoplus_{g\in G}D .
\end{equation}
Consequently $\Msym^G\cong D$ is a factor,
\begin{equation}\label{eq:index}
\bigl[\Msym:\Msym^{G}\bigr]_0=|G|=|\Xi_\beta|,
\qquad
\Msym^{G}{}'\cap\Msym=Z(\Msym)\cong\C^{|G|},
\end{equation}
the index being realized by $E(x)=|G|^{-1}\sum_{g}\theta_g(x)$; and the intermediate von
Neumann algebras are exactly the $D\otimes\ell^\infty(G/\mathcal P)$, $\mathcal P$ a
partition of $G$, so $\Lat(\Msym^G\subset\Msym)$ is the partition lattice of $G$. The
subgroup lattice of $G$ embeds in it, order-reversingly, by $H\mapsto\Msym^H$, which
corresponds to the partition of $G$ into cosets of $H$; for $|G|\ge3$ the embedding is
proper.
\end{theorem}

\begin{proof}
$G$ permutes the summands of \eqref{eq:central} simply transitively, so an invariant element
is determined by one component and $\Msym^G\cong D$ embedded diagonally via the $\theta_g$,
which is \eqref{eq:trivialincl}. In \eqref{eq:trivialincl} the averaging map is
$D\otimes E_0$ with $E_0$ the uniform state on $\ell^\infty(G)$, whose Pimsner--Popa index
is $|G|$; any other expectation is $D\otimes E_\mu$ with $\mu$ a faithful probability on
$G$ and index $\max_g\mu(g)^{-1}\ge|G|$, so the Kosaki--Longo index is $|G|$. An element of
$\Msym$ commuting with the diagonal copy of the factor $D$ is scalar in each summand, giving
$\C^{|G|}$. By the Ge--Kadison splitting theorem \cite{GK}, every von Neumann algebra
between $D\otimes1$ and $D\otimes\ell^\infty(G)$ is $D\otimes B$ for a von Neumann
subalgebra $B\subseteq\ell^\infty(G)$, and the unital subalgebras of $\ell^\infty(G)$ are the
$\ell^\infty(G/\mathcal P)$. The $H$-invariant elements are those constant on the
$H$-orbits, i.e.\ on the cosets of $H$; the partition $\{1\},\ G\setminus\{1\}$ of a group of
order $\ge3$ is not a coset partition.
\end{proof}

\begin{remark}[Not a subfactor invariant]\label{rem:nostandard}
The inclusion \eqref{eq:trivialincl} is the trivial inclusion of the factor $D$ tensored
with $\C\subseteq\ell^\infty(G)$: its Jones tower is $D\otimes(\C\subseteq\ell^\infty(G)\subseteq
M_{|G|}(\C)\subseteq\cdots)$, and its intermediate algebras are governed by partitions rather
than by subgroups. It is therefore \emph{not} of the form $M^G\subset M$ for an outer action
on a factor, and it carries no fusion category of bimodules beyond the data of the finite set
$G$. The arithmetic content of the symmetry breaking is entirely in the number $|\Xi_\beta|$
and in the way $G_S$ permutes the summands; this is what Corollary~\ref{cor:staircase}
quantifies. Along a cascade with $\Xi_\beta\cong(\Z/2\Z)^{A_\beta}$ \cite[Prop.~4.16]{Main},
crossing a threshold $\beta_n$ doubles the number of summands and doubles the index.
\end{remark}

\begin{lemma}[Relative entropy of the phases]\label{lem:relent}
With the notation of Lemma~\ref{lem:central}, and $\Ent$ the relative entropy,
\[
\Ent\bigl(\omega_{\varphi_\varepsilon}\,\big\|\,\omega_{\phisym}\bigr)=\log|G| ,
\]
and the whole value is carried by the centre: the contribution of each summand
$M_\varepsilon$ vanishes. For $\varepsilon\ne\varepsilon'$ one has
$\Ent(\omega_{\varphi_\varepsilon}\|\omega_{\varphi_{\varepsilon'}})=+\infty$.
\end{lemma}

\begin{proof}
For a direct sum $M=\bigoplus_iM_i$ with normal states $\omega=\bigoplus\lambda_i\omega_i$
and $\psi=\bigoplus\mu_i\psi_i$,
\[
\Ent(\omega\|\psi)=\sum_i\lambda_i\log\frac{\lambda_i}{\mu_i}+\sum_i\lambda_i\,\Ent(\omega_i\|\psi_i),
\]
with the conventions $0\log0=0$ and $\lambda\log(\lambda/0)=+\infty$ for $\lambda>0$. Apply
this to \eqref{eq:central} with $\lambda=\delta_\varepsilon$ and $\mu_\delta=|G|^{-1}$ for
every $\delta$, which is Lemma~\ref{lem:central}. On the summand $\varepsilon$ the two
states induce the \emph{same} normal state of $M_\varepsilon$, so that term vanishes; on the
others $\lambda_\delta=0$. Hence $\Ent=\log|G|$ and the second sum is identically zero,
which is the assertion that the relative entropy is purely central. For
$\varepsilon\ne\varepsilon'$ the two states are supported on orthogonal central projections
and the first sum is $\log(1/0)=+\infty$.
\end{proof}

\begin{corollary}[Index staircase equals information staircase]\label{cor:staircase}
\[
\log\bigl[\Msym:\Msym^G\bigr]_0=|A_\beta|\log2=\Ent\bigl(\omega_{\varphi_\varepsilon}\big\|\,\omega_{\phisym}\bigr),
\]
by \eqref{eq:index} and Lemma~\ref{lem:relent}. The subfactor index and the relative
entropy are the same staircase.

The coincidence is not an accident, but neither is it deep: both sides count phases. The
index of $\Msym^G\subset\Msym$ is the number $|G|$ of extremal $\mathrm{KMS}_\beta$ states,
by \eqref{eq:trivialincl}; and the relative entropy of one extremal state against the
uniform mixture of $|G|$ mutually singular states is $\log|G|$, the information needed to
name which phase one is in. Each threshold $\beta_n$ doubles the number of phases, so each
adds $\log2$ to both quantities. What the identity does say is that the operator-algebraic
and the information-theoretic bookkeeping of the symmetry breaking agree exactly, with no
correction term from the type $\mathrm{III}$ structure --- a point made independently in
Lemma~\ref{lem:relent}, where the relative entropy is shown to be purely central.
\end{corollary}

\section{The control experiment: restoring the additive part}\label{sec:affine}

To confirm that Lemma~\ref{lem:single} is the cause of \S\S2--5, we restore the missing
ingredient.

\begin{remark}[The additive group is $\cO_K$]\label{rem:additive}
The ring of $S$-integers $\cO_{K,S}$ inverts the primes of $S$, so its elements have poles
at $S$ and it is not contained in $\hO_S$. The correct additive part is $\cO_K$, dense in
$\hO_S$ by the Chinese remainder theorem. We let $J_S$ act on $\hO_S$ through a fixed choice
of uniformizers $\pi_\pp\in K_\pp$, $\aaa=\prod\pp^{a_\pp}$ acting by
$x\mapsto\pi_\aaa x$ with $(\pi_\aaa x)_\pp=\pi_\pp^{a_\pp}x_\pp$, and set
$\cA^{\mathrm{aff}}_{K,S}=C(\hO_S)\rtimes(\cO_K\rtimes J_S)$, the semigroup crossed product
by the maps $x\mapsto\pi_\aaa x+b$. Over $\Q$ the choice $\pi_\ell=\ell$ is canonical; in
general everything below holds for every choice. (Without the additive part, the choice is
absorbed by the $G_S$-coordinate of $Y_S$; with it, the algebra is that of Cuntz--Li
\cite{CL} and Cuntz--Deninger--Laca \cite{CDL} localized at $S$.)
\end{remark}

\begin{lemma}[Finite covers]\label{lem:cover}
For every $\pp\in S$, $\hO_S=\bigsqcup_{a\in\cO_K/\pp}(a+\pp\hO_S)$, a partition into
$\Ns\pp$ clopen sets. The isometries $\mu_{(a,\pp)}$ therefore have mutually orthogonal
ranges and satisfy the Cuntz relation
\begin{equation}\label{eq:cuntz}
\sum_{a\in\cO_K/\pp}\mu_{(a,\pp)}\mu_{(a,\pp)}^*=1 ,
\end{equation}
and the tight spectrum is all of $\hO_S$: no collapse occurs, and
$\cA^{\mathrm{aff}}_{K,S}$ is its own boundary quotient, $\partial\cA^{\mathrm{aff}}_{K,S}=\cA^{\mathrm{aff}}_{K,S}$.
\end{lemma}

\begin{proof}
$\pp\hO_S$ is open of index $\Ns\pp$ and $\cO_K$ surjects onto $\hO_S/\pp\hO_S$ by density.
The cover is finite, so the tight relations of Exel \cite{Exel} impose equality rather
than the nested inequality of Lemma~\ref{lem:single}; and \eqref{eq:cuntz} already holds in
$\cA^{\mathrm{aff}}_{K,S}$, since the Stacey covariance gives
$\mu_{(a,\pp)}\mu_{(a,\pp)}^*=\mathbf 1_{a+\pp\hO_S}$.
\end{proof}

\begin{theorem}[Kirchberg boundary]\label{thm:kirchberg}
The action of $\cO_K\rtimes J_S$ on $\hO_S$ is minimal and topologically free, and
$\partial\cA^{\mathrm{aff}}_{K,S}$ is a unital separable nuclear simple purely infinite
$C^*$-algebra satisfying the UCT.
\end{theorem}

\begin{proof}
Minimality: $\cO_K$ is dense and acts by translations, so every orbit is dense.
Topological freeness: a nontrivial $(b,\aaa)$ with $\aaa=1$ is a nontrivial translation of a
group, hence fixed-point free. For $\aaa\ne1$ the fixed points satisfy $(1-\pi_\aaa)x=b$
componentwise: at $\pp\in\operatorname{supp}\aaa$ the element $1-\pi_\pp^{a_\pp}$ is a unit,
so $x_\pp$ is determined; at $\pp\notin\operatorname{supp}\aaa$ the equation reads $0=b$ in
$\cO_\pp$, which forces $b=0$. So the fixed-point set is empty, or a single point, or (for
$b=0$) the set $\{x:x_\pp=0\ \forall\pp\in\operatorname{supp}\aaa\}$; in every case it is closed
with empty interior. The union over the countable semigroup is meagre, hence has empty
interior. The groupoid of germs is \'etale, Hausdorff and amenable, being that of a partial
action of the amenable group $K\rtimes I_S$; minimality and topological freeness give
simplicity \cite{BCFS}. It is locally contracting: a basic clopen set
$V=\{x:x_\qq\in a_\qq+\qq^{n_\qq}\cO_\qq,\ \qq\in F\}$ is mapped by $x\mapsto\pi_\pp x+c$,
$\pp\in F$ and $c\in\cO_K$ chosen by the Chinese remainder theorem with
$c\equiv(1-\pi_\pp)a_\pp\pmod{\pp^{n_\pp}}$ and $c\equiv0\pmod{\qq^{n_\qq}}$ for $\qq\ne\pp$,
onto a clopen proper subset of $V$. Hence $\cA^{\mathrm{aff}}_{K,S}$ is purely infinite by
Anantharaman-Delaroche \cite{AD}. Nuclearity and the UCT follow from amenability.
\end{proof}

\begin{theorem}[The unit class says little even here]\label{thm:unit}
In $K_0(\partial\cA^{\mathrm{aff}}_{K,S})$ one has $(\Ns\pp-1)[1]=0$ for every $\pp\in S$,
so the order of $[1]$ divides $d_S=\gcd_{\pp\in S}(\Ns\pp-1)$. In particular $[1]=0$ whenever
$d_S=1$, which holds as soon as $S$ contains a prime of norm $2$, or two primes with
$\gcd(\Ns\pp-1,\Ns\qq-1)=1$.
\end{theorem}

\begin{proof}
Each $\mu_{(a,\pp)}$ is an isometry, so $[\mu\mu^*]=[1]$; applying $K_0$ to \eqref{eq:cuntz}
gives $[1]=\Ns\pp[1]$.
\end{proof}

\begin{proposition}[The Cuntz relation pins the temperature]\label{prop:betaone}
For the dynamics $\sigma_t(\mu_{(b,\aaa)})=\Ns\aaa^{it}\mu_{(b,\aaa)}$ on
$\cA^{\mathrm{aff}}_{K,S}$, a $\mathrm{KMS}_\beta$ state exists only for $\beta=1$.
\end{proposition}

\begin{proof}
For a $\mathrm{KMS}_\beta$ state $\varphi$ and any of the $\Ns\pp$ isometries in
\eqref{eq:cuntz}, the KMS condition gives
$\varphi(\mu\mu^*)=\varphi(\mu^*\sigma_{i\beta}(\mu))=\Ns\pp^{-\beta}\varphi(\mu^*\mu)=\Ns\pp^{-\beta}$;
summing \eqref{eq:cuntz} gives $1=\Ns\pp^{1-\beta}$, so $\beta=1$.
\end{proof}

\begin{example}
Over $\Q$: $S=\{3\}$ gives $d_S=2$; $S=\{3,5\}$ gives $\gcd(2,4)=2$; $S=\{7,13\}$ gives
$\gcd(6,12)=6$; but any $S$ containing $2$ gives $d_S=1$ and $[1]=0$. A single small prime
erases the invariant.
\end{example}

\begin{remark}[Two behaviours of the unit class]\label{rem:twounits}
Theorems \ref{thm:Kboundary} and \ref{thm:unit} exhibit opposite behaviours of $[1]$, and
the contrast is instructive.

In the multiplicative boundary, $K_0=\Lambda^{\mathrm{ev}}(I_S^0)$ is torsion-free and
$[1]=h^+_S\,e_0$ is a nonzero element of divisibility exactly $h^+_S$: the narrow class
number is stored in \emph{how divisible} the unit is. In the affine boundary the Cuntz
relation forces $[1]=\Ns\pp\,[1]$, so $[1]$ is a torsion element of order dividing
$\gcd_\pp(\Ns\pp-1)$ and is usually zero. The mechanism is already visible for the Cuntz
algebra $\cO_n$, where $K_0=\Z/(n-1)$ is generated by $[1]$; taking several primes replaces
$n-1$ by a greatest common divisor, and a single prime of norm $2$ collapses it.

So the two systems store arithmetic in $[1]$ in incompatible ways, and only the
multiplicative one does so usefully. This is a rare instance in which the nested range
projections of Lemma~\ref{lem:single} are an advantage: they leave $K_0$ torsion-free, and a
torsion-free group can record an integer by divisibility, which a group killed by
$\gcd_\pp(\Ns\pp-1)$ cannot. We have not computed the ambient groups
$K_*(\partial\cA^{\mathrm{aff}}_{K,S})$ and make no claim about them beyond
Theorem~\ref{thm:unit}.
\end{remark}

\section{Assessment}

\[
\begin{array}{lll}
\text{Feature} & \cA_{K,S}\ (\text{multiplicative}) & \cA^{\mathrm{aff}}_{K,S}\ (\text{affine})\\\hline
\text{isometries per }\pp & 1 & \Ns\pp\\
\text{range projections} & \text{nested} & \text{partitioning}\\
\text{tight relation} & \mu_\pp\ \text{unitary} & \text{Cuntz}\\
\text{boundary space} & \Cl^+_S,\ \text{finite} & \hO_S,\ \text{a Cantor group}\\
\text{boundary algebra} & M_{h^+_S}(C^*(I_S^0)),\ \text{type }\mathrm I & \text{Kirchberg (the algebra itself)}\\
\text{trace on the boundary} & \text{faithful} & \text{none}\\
{[1]}\ \text{records} & h^+_S & \text{order}\mid\gcd(\Ns\pp-1)
\end{array}
\]

Everything in the left column follows from Lemma~\ref{lem:single}, and the right column
confirms by contrast that the additive cosets are the cause. The same dichotomy governs the
admissible inverse temperatures --- the Cuntz relation forces $\beta=1$ on the affine
system (Proposition~\ref{prop:betaone}), whereas the multiplicative system has
$\mathrm{KMS}_\beta$ states for every $\beta>0$ \cite[Thm.~3.6]{Main} --- and the free
variant treated in Paper V of the series, where orthogonal ranges again restrict $\beta$.

Four questions remain. We have computed $K_*(\cA_{K,S})$ only for a single prime, and
$K_*(\cA^{\mathrm{aff}}_{K,S})$ not at all; by Theorem~\ref{thm:kirchberg} the triple
$(K_0,[1],K_1)$ is a complete invariant of the affine algebra and by Theorem~\ref{thm:unit}
the middle entry is usually trivial, so the content lies in the groups. When
$\cO^\times_{K,+}$ is infinite, whether some proper stratum carries nontrivial isotropy
(Remark~\ref{rem:leopoldt}) --- a Leopoldt-type closure question --- and hence contributes
further ideals of the kind seen in Theorem~\ref{thm:principal}(2), is open. Whether the
infinite index of Theorem~\ref{thm:infinite} persists at the non-symmetric extremal states
(Remark~\ref{rem:scope}), and whether these inclusions are discrete in the sense of
Izumi--Longo--Popa \cite{ILP} or admit a continuous decomposition matching the diffuse
relative commutant, is open. Finally, the lattice of intermediate subfactors
\cite{Watatani} of the symmetry-breaking inclusion is the partition lattice of
$\widehat{\Xi_\beta}$ (Theorem~\ref{thm:finite}), which remembers $\Xi_\beta$ only as a finite
set; it distinguishes neither $K$ nor $S$ nor the thresholds, and the question of which
invariant does is taken up in Paper III of the series.

It is worth saying plainly why the blindness recorded in Remark~\ref{rem:noarith} is not a
computational accident. By Propositions \ref{prop:K} and \ref{prop:Kfinite} the passage
from $C(Y_S)$ to $\cA_{K,S}$ is, on the image of $K^0(Y_S)$ in $K_0$, the passage to the
coinvariants of $\alpha_*$, and a functional survives that quotient exactly when it is
$\alpha_*$-invariant. The arithmetically
meaningful functionals are the $\mathrm{KMS}_\beta$ measures, which by
\cite[Prop.~2.8]{Main} are eigenvectors,
$\omega\circ\alpha_{\pp*}=\Ns\pp^{-\beta}\omega$, with eigenvalue $\ne1$ for $\beta\ne0$.
Passing to the coinvariants therefore \emph{annihilates precisely the covariant vectors that
carry the arithmetic}, by construction rather than by accident. The same mechanism in its
measure-theoretic form is, we expect, the reason $\Xi_\beta$ is an
invariant of a measure class rather than of a $C^*$-algebra: any functor that forgets
$\alpha_*$, or that forces it to act as the identity, loses it.

\begin{thebibliography}{9}
\bibitem{AD} C. Anantharaman-Delaroche, \emph{Purely infinite $C^*$-algebras arising from dynamical systems}, Bull. Soc. Math. France 125 (1997), 199--225.
\bibitem{BCFS} J. Brown, L. O. Clark, C. Farthing, A. Sims, \emph{Simplicity of algebras associated to \'etale groupoids}, Semigroup Forum 88 (2014), 433--452.
\bibitem{CDL} J. Cuntz, C. Deninger, M. Laca, \emph{$C^*$-algebras of Toeplitz type associated with algebraic number fields}, Math. Ann. 355 (2013), 1383--1423.
\bibitem{CL} J. Cuntz, X. Li, \emph{The regular $C^*$-algebra of an integral domain}, in: Quanta of Maths, Clay Math. Proc. 11, AMS, 2010, 149--170.
\bibitem{Exel} R. Exel, \emph{Inverse semigroups and combinatorial $C^*$-algebras}, Bull. Braz. Math. Soc. 39 (2008), 191--313.
\bibitem{GK} L. Ge, R. Kadison, \emph{On tensor products of von Neumann algebras}, Invent. Math. 123 (1996), 453--466.
\bibitem{Green} P. Green, \emph{The local structure of twisted covariance algebras}, Acta Math. 140 (1978), 191--250.
\bibitem{ILP} M. Izumi, R. Longo, S. Popa, \emph{A Galois correspondence for compact groups of automorphisms of von Neumann algebras}, J. Funct. Anal. 155 (1998), 25--63.
\bibitem{Kasparov} G. G. Kasparov, \emph{Equivariant $KK$-theory and the Novikov conjecture}, Invent. Math. 91 (1988), 147--201.
\bibitem{Katsura} T. Katsura, \emph{On $C^*$-algebras associated with $C^*$-correspondences}, J. Funct. Anal. 217 (2004), 366--401.
\bibitem{Kosaki} H. Kosaki, \emph{Extension of Jones' theory on index to arbitrary factors}, J. Funct. Anal. 66 (1986), 123--140.
\bibitem{Laca} M. Laca, \emph{From endomorphisms to automorphisms and back: dilations and full corners}, J. London Math. Soc. (2) 61 (2000), 893--904.
\bibitem{Longo} R. Longo, \emph{Index of subfactors and statistics of quantum fields I}, Comm. Math. Phys. 126 (1989), 217--247.
\bibitem{Pimsner} M. V. Pimsner, \emph{A class of $C^*$-algebras generalizing both Cuntz--Krieger algebras and crossed products by $\Z$}, in: Free Probability Theory, Fields Inst. Commun. 12, AMS, 1997, 189--212.
\bibitem{PV} M. Pimsner, D. Voiculescu, \emph{Exact sequences for $K$-groups and $\mathrm{Ext}$-groups of certain cross-product $C^*$-algebras}, J. Operator Theory 4 (1980), no.~1, 93--118.
\bibitem{PP} M. Pimsner, S. Popa, \emph{Entropy and index for subfactors}, Ann. Sci. \'Ec. Norm. Sup. (4) 19 (1986), 57--106.
\bibitem{Renault} J. Renault, \emph{The ideal structure of groupoid crossed product $C^*$-algebras}, J. Operator Theory 25 (1991), 3--36.
\bibitem{Renault08} J. Renault, \emph{Cartan subalgebras in $C^*$-algebras}, Irish Math. Soc. Bull. 61 (2008), 29--63.
\bibitem{Takesaki} M. Takesaki, \emph{Conditional expectations in von Neumann algebras}, J. Funct. Anal. 9 (1972), 306--321.
\bibitem{Watatani} Y. Watatani, \emph{Lattices of intermediate subfactors}, J. Funct. Anal. 140 (1996), 312--334.
\bibitem{Main} R. Chen, \emph{KMS state uniqueness and phase transitions in localized Bost--Connes systems: from Chebotarev density to the Hardy--Littlewood conjecture}, preprint, 2026, DOI \href{https://doi.org/10.5281/zenodo.22152101}{10.5281/zenodo.22152101}.
\end{thebibliography}

\end{document}
